% ===================================================================
% Chapter 5: Normal contributions and Table III, Eqs. (I.68)-(I.93)
% Tasks 11 and 12 of the re-derivation.
% ===================================================================
\providecommand{\chVdd}{\nabla\!\cdot\!\nabla'}          % grad . grad'
\providecommand{\chVlead}{A_{\rm lead}}                  % leading term of A(y)
\providecommand{\chVAh}{\hat A}                          % A(y)/h at D=4
\providecommand{\chVyes}{\textcolor{tagD}{agrees}}
\providecommand{\chVno}{\textcolor{tagC}{differs}}

\chapter{Normal contributions and Table III: Eqs. (I.68)--(I.93)}
\label{ch:05}
% Equations tagged (I.nn) do not step the equation counter; with hyperref's default anchor
% names they would share anchors with the next numbered equation. Use unique anchors here.
\hypersetup{hypertexnames=false}

This chapter covers the first part of the ``normal'' 3-point contributions of \PI. These are the terms
\eqI{53} in which two derivatives act on one propagator function and the result is multiplied by a second
propagator function $f(y)$. \PI{} reduces them by extracting the derivatives \emph{generically}, using the
indefinite integral $I[f]$ of \eqI{68}.
\begin{itemize}[leftmargin=14pt]
\item \textbf{Task 11} (Section~\ref{sec:ch05:t11}) re-derives the derivative-extraction identities \eqI{68}--\eqI{80}.
It also examines in detail the statement \eqI{69} that no delta function is produced.
\item \textbf{Task 12} (Section~\ref{sec:ch05:t12}) re-derives the time-derivative identities \eqI{81}--\eqI{88}.
It applies them to the nonlocal logarithm term of \eqI{67}, which gives \eqI{89} and Table~III.
It then derives \eqI{90}--\eqI{93}, which prepare the reduction of the mixed $A$ terms; the reduced normal contributions are collected in Tables~IV and V (next chapter).
\end{itemize}
The identities \eqI{68}--\eqI{88} and \eqI{90}--\eqI{93} do not depend on any sign convention.
Only \eqI{89} and Table~III carry the overall sign of the self-mass. For them we state explicitly whether an equation
is in the \emph{paper convention} or in the \emph{definition sign} fixed by \eqI{168}.

\paragraph{Notation used in this chapter.}
$\Delta x^\mu\equiv x^\mu-x'^\mu$. The lower-index form is $\Delta x_\sigma\equiv\eta_{\sigma\nu}\Delta x^\nu$, so that $\Delta x_0=-\Delta\eta$, with $\Delta\eta\equiv\eta-\eta'$.
We write $\|\Delta\vec x\|^2=\sum_i(\Delta x^i)^2$. The functions $f$, $A$, $F$, $G$, $K$ depend on $y$ only, and primes on them denote $d/dy$.
$\partial$ and $\partial'$ are derivatives with respect to $x$ and $x'$. Since $a=-1/(H\eta)$, we have $\partial_0a=Ha^2$.
The constants are those of earlier chapters:
\begin{equation*}
\begin{gathered}
\mathcal N\equiv\frac{H^{D-2}}{(4\pi)^{D/2}},\qquad
k\equiv\mathcal N\,\frac{\Gamma(D-1)}{\Gamma(\frac D2)}\xrightarrow{D\to4}\frac{H^2}{8\pi^2},\\
\pihat\equiv-\pi\cot\big(\tfrac{\pi D}{2}\big)=-\frac{2}{D-4}+\order{D-4},\qquad
h\equiv\frac{H^2}{16\pi^2}.
\end{gathered}
\end{equation*}
The leading term of the massless minimally coupled propagator function is
$\chVlead\equiv\mathcal N\,\Gamma(\tfrac D2-1)\,(4/y)^{\frac D2-1}$. Here the first coefficient of \eqI{14} is the corrected value
$\Gamma(\frac D2-1)$ established in Task~8; the paper prints $\Gamma(\frac D2-1)/(\frac D2-1)$, a typographical error.
$I[f]$ is the indefinite integral \eqI{68}, and $I^n$ is its $n$-fold iterate.
Results of earlier tasks are cited by task number together with the equation or section of these notes where they are derived: Task~4 in Chapter~\ref{ch:fourpoint}, Tasks~7 and~8 in Chapter~\ref{ch:three-point-setup}, and Tasks~9 and~10 in Chapter~\ref{ch:ch04}.

% ===================================================================
\section[Task 11: derivative-extraction identities, Eqs. (I.68)--(I.80)]{Task 11: derivative-extraction identities for the normal contributions, Eqs.~(I.68)--(I.80)}
\label{sec:ch05:t11}

\begin{resultbox}[Summary of findings (Task 11)]
\begin{itemize}[leftmargin=12pt,itemsep=2pt]
\item \textbf{\eqI{68}--\eqI{80}} are identities that do not depend on the sign convention: neither the factor $\kappa^2$ nor the sign question of \eqI{47} enters.
All of them were derived and checked. \eqI{71}--\eqI{73} also agree term by term with the derivative table \eqref{eq:ch3-yder} of Task~8.
\item \textbf{A qualification of \eqI{69}.} The statement that two derivatives acting on a product of two propagators never produce a delta function holds for the \emph{naive} (pointwise) derivative.
In dimensional regularization the \emph{distributional} derivative produces a term
$\delta^0_\rho\delta^0_\sigma\,\delta^D(x-x')$ from the component of $F(y)$ with exactly the power $y^{1-D/2}$.
For $I^2[fA'']$ this component is $c_f\,\chVlead$, where $c_f$ is the constant part of $f$.
The distributional delta content is therefore \emph{exactly the same} as the local contribution counted separately in \eqI{54}--\eqI{58}.
Consequently \eqI{70} holds in both readings, provided both sides are naive or both are distributional.
But since the local contribution is counted separately, the derivatives in the normal contributions must be taken naively.
Otherwise the local term, which has a pole for $I=A$, would be counted twice.
Whether the paper follows this rule in Tables~IV--VI is to be checked there.
\item \textbf{The divergence count after \eqI{80} is correct.} At $D=4$ only $y^{2-D}$ in $I^2[AA'']$, and only $y^{1-D}$ and $y^{2-D}$ in $I[A'^2]$, fail to be locally integrable.
In all other terms one may set $D=4$.
\end{itemize}
\end{resultbox}

% -------------------------------------------------------------------
\subsection{Goal, inputs and scope}

\paragraph{Targets \tP.}
The targets are Eqs.~\eqI{68}--\eqI{69} (p.~124005-7, left column), \eqI{70}--\eqI{78} (p.~124005-7, right column) and \eqI{79}--\eqI{80} (p.~124005-8, left column).
A footnote to \eqI{75} specifies that its left-hand side means the naive second derivative, without the delta function.
The printed equations were \emph{not} used as input to the derivation. They enter only in the comparison, Section~\ref{sec:ch05:t11:comp}.
\eqI{68} is a definition.

\paragraph{Inputs.}
\begin{itemize}[leftmargin=14pt]
\item[(G1)] \tP\ \eqI{11}: $y=aa'H^2\big\{\|\Delta\vec x\|^2-(|\Delta\eta|-i\delta)^2\big\}$ with $a=-1/(H\eta)$.
\item[(G2)] \tD\ \eqI{52} and its split into a local part and a regular part (Task~8). In particular, at fixed $\delta>0$ the singular part is
\begin{equation}
\Big[\partial_0\partial'_0F(y)\Big]_{\rm sing}=-4i\delta\,aa'H^2\,F'(y)\Big|_{\Delta\eta=0}\,\delta(\Delta\eta)
\qquad\text{(Task 8, Eqs.\ \eqref{eq:ch3-chain}--\eqref{eq:ch3-Sdelta})}.
\label{eq:ch05:t11:sing}
\end{equation}
\item[(G3)] \tP\ The definition \eqI{53} of the normal contribution: $f(y)\times$[the regular part of \eqI{52}], with $f\in\{A,B,C\}$.
For $I=A$ we have $f=A(y)$. The logarithm part is treated separately (Task~10).
\end{itemize}

\paragraph{Scope.}
\eqI{81} and beyond, and Tables~III--V, are not treated in this task. All identities are derived in general $D$.

% -------------------------------------------------------------------
\subsection[Eq. (I.68): the indefinite integral]{Eq.~(I.68): the indefinite integral \tA}

\begin{equation}
I[f](y)\equiv\int^y dy'\,f(y'),\qquad \frac{d}{dy}I[f]=f,\qquad I^n\equiv I\circ\cdots\circ I .
\tag{I.68}\label{eq:ch05:I68}
\end{equation}
\paragraph{Ambiguity of the constants.}
$I[f]$ is defined only up to a constant, and $I^2[f]$ only up to $c_1y+c_0$. In each identity below this ambiguity cancels when the constants are chosen consistently.
For example, in \eqI{70} the change $I^2[fA'']\to I^2[fA'']+c_1y+c_0$ produces $c_1\,\partial_\rho\partial'_\sigma y$ under $\partial_\rho\partial'_\sigma$.
This is compensated by $I[f'A']\to I[f'A']-c_1$. The results are therefore independent of the integration constants.

% -------------------------------------------------------------------
\subsection[Eq. (I.69): chain rule and delta-function content]{Eq.~(I.69): chain rule and delta-function content \tD}
\label{sec:ch05:t11:69}

\paragraph{Naive derivative.}
As long as $F$ is smooth, the chain rule holds pointwise:
\begin{equation}
\partial_\rho\partial'_\sigma F(y)=F''(y)\,\frac{\partial y}{\partial x^\rho}\frac{\partial y}{\partial x'^\sigma}+F'(y)\,\frac{\partial^2y}{\partial x^\rho\partial x'^\sigma}.
\tag{I.69}\label{eq:ch05:I69}
\end{equation}

\paragraph{Distributional derivative.}
As in Task~8, at fixed $\delta>0$ a single singular term comes from the kink of $|\Delta\eta|$:
$\delta^0_\rho\delta^0_\sigma\,(-4i\delta\,a^2H^2)\,F'(y)\big|_{\Delta\eta=0}\,\delta(\Delta\eta)$.
At $\Delta\eta=0$ one has $y=a^2H^2(r^2+\delta^2)$ with $r=\|\Delta\vec x\|$.
For a component $F\supset y^s$, so that $F'\propto y^{s-1}$, the $d^{D-1}x$ integral of this term scales as
\begin{equation}
\delta\int d^{D-1}x\,(r^2+\delta^2)^{s-1}\;\propto\;\delta^{\,D+2s-2}.
\label{eq:ch05:t11:scaling}
\end{equation}
Consequently:
\begin{itemize}[leftmargin=14pt]
\item If $s=1-\frac D2$, the exponent vanishes and a finite $\delta^D$ term remains (Task~8, Eq.~\eqref{eq:ch3-lead}).
\item For every other $s$, the exponent $D+2s-2$ is a nonzero analytic function of $D$. In the region where its real part is positive the limit $\delta\to0$ is zero.
In dimensional regularization it is then zero for general $D$ by analytic continuation \tA.
\end{itemize}
Hence
\begin{equation}
\Big[\partial_\rho\partial'_\sigma F\Big]_{\rm distr}=\eqref{eq:ch05:I69}+\delta^0_\rho\delta^0_\sigma\,\gamma_F\,\frac{i}{a^{D-2}}\,\delta^D(x-x'),
\qquad F\supset\gamma_F\cdot\mathcal N\,\Gamma(\tfrac D2-1)\Big(\frac4y\Big)^{\frac D2-1}.
\label{eq:ch05:t11:distr}
\end{equation}
The normalization of $\gamma_F$ is chosen so that $F=\chVlead$ gives $\gamma_F=1$, with the local coefficient $i\,a^{2-D}$ fixed in Task~8.

\paragraph{Conclusion.}
\eqI{69} holds for the naive derivative. For the distributional derivative it holds only if $F$ has no $y^{1-D/2}$ component.
A product of two propagators $i\Delta_I\,i\Delta_A$ \emph{does} contain this component: the constant part of one factor times the leading term of the other.
So $\gamma_F\neq0$. Section~\ref{sec:ch05:t11:delta} shows why this matters.

% -------------------------------------------------------------------
\subsection[Eq. (I.70)]{Eq.~(I.70) \tD}

Take $G\equiv I^2[fA'']$, so that $G'=I[fA'']$ and $G''=fA''$. By \eqref{eq:ch05:I69},
\begin{equation*}
\partial_\rho\partial'_\sigma I^2[fA'']=fA''\,\partial_\rho y\,\partial'_\sigma y+I[fA'']\,\partial_\rho\partial'_\sigma y .
\end{equation*}
Integration by parts gives $I[fA'']=fA'-I[f'A']$, with consistent constants. Substituting:
\begin{resultbox}
\begin{equation}
f(y)\Big\{A''(y)\,\frac{\partial y}{\partial x^\rho}\frac{\partial y}{\partial x'^\sigma}+A'(y)\,\frac{\partial^2y}{\partial x^\rho\partial x'^\sigma}\Big\}
=\partial_\rho\partial'_\sigma I^2[fA''](y)+\frac{\partial^2y}{\partial x^\rho\partial x'^\sigma}\,I[f'A'](y).
\tag{I.70}\label{eq:ch05:I70}
\end{equation}
\end{resultbox}
\noindent This is Eq.~\eqI{70} of the paper, identical to the printed equation.

% -------------------------------------------------------------------
\subsection[Delta content of (I.70) and consistency with the local contribution]{Delta content of (I.70) and consistency with the local contribution \tD}
\label{sec:ch05:t11:delta}

Write $f=c_f+(\text{powers of }y)$, where $c_f$ is the constant part of $f$.
From Task~4, $c_f=k\pihat$ for $I=A$, $c_f=b$ for $I=B$ and $c_f=c$ for $I=C$.
In general $D$, the $D$-dependent powers do not mix with the integer powers. Therefore the $y^{1-D/2}$ component of $I^2[fA'']$ comes only from $c_f\times\chVlead''$:
\begin{equation}
\begin{aligned}
&I^2[fA'']\supset I^2\big[c_f\,\chVlead''\big]=c_f\,\chVlead\ \ (+\text{linear in }y),\qquad \chVlead=\mathcal N\,\Gamma(\tfrac D2-1)\,(4/y)^{\frac D2-1}\\
&\Longrightarrow\quad\gamma_{I^2[fA'']}=c_f .
\end{aligned}
\label{eq:ch05:t11:gamma}
\end{equation}
Moreover, $I[f'A']$ contains no $y^{-D/2}$ component. Its product with the $\order{\delta}$ singular part of $\partial_0\partial'_0y$ therefore gives no delta function either.
(Editorial gloss: by the scaling \eqref{eq:ch05:t11:scaling}, a finite $\delta^D$ from $\delta\cdot y^{s'}$ requires $s'=-D/2$. The powers present in $I[f'A']$ are non-negative integer powers, $y^{n-D/2}$ with $n\geq1$, and $y^{m-D}$ with $m\geq1$.)
From this:
\begin{itemize}[leftmargin=14pt]
\item \textbf{Naive reading.} Both sides of \eqref{eq:ch05:I70} hold with naive derivatives. The left-hand side is the normal contribution of \eqI{53}.
\item \textbf{Distributional reading.} If the full $\partial\partial'A$ is taken on the left, a term $f(y)\cdot\frac{i}{a^{D-2}}\delta^D=c_f\,\frac{i}{a^{D-2}}\delta^D$ is added.
This uses $f(y)\delta^D=f(0)\delta^D=c_f\delta^D$, which holds in dimensional regularization.
By \eqref{eq:ch05:t11:distr} and \eqref{eq:ch05:t11:gamma}, exactly the same term is added on the right. So \eqref{eq:ch05:I70} also holds in the distributional sense.
\end{itemize}
\begin{discbox}[Conclusion: naive inner derivatives in the normal contributions]
The distributional $\partial\partial'I^2[fA'']$ contains the \emph{whole} local contribution \eqI{54}--\eqI{58} of Task~9.
For $I=A$ this is \eqI{54}, which carries the $k\pihat$ pole.
Since the local contribution has been computed separately, $\partial\partial'I^2[fA'']$ must be taken naively in the normal contributions.
Suppose that in later steps $\square$ were applied as a distribution to the $c_f\chVlead$ part of $I^2[fA'']$ (for example by using the fact that $\square(4/y)^{D/2-1}$ produces a $\delta^D$ term).
The local term would then be counted twice.
The statement ``never produces a delta function'' in \eqI{69} is consistent with this naive reading.
When Tables~IV--VI are derived, it must be checked whether this rule is respected.
\end{discbox}

% -------------------------------------------------------------------
\subsection[Eqs. (I.71)--(I.73): the mixed second derivative of y]{Eqs.~(I.71)--(I.73): the mixed second derivative of $y$ \tD}

Take the regular part ($\delta\to0$), $y=aa'H^2\big(\|\Delta\vec x\|^2-\Delta\eta^2\big)$, and use $\partial_0a=Ha^2$:
\begin{equation}
\partial_0y=Ha\,y+2aa'H^2\Delta x_0,\qquad \partial'_0y=Ha'\,y-2aa'H^2\Delta x_0,\qquad \partial_iy=2aa'H^2\Delta x_i=-\partial'_iy .
\label{eq:ch05:t11:dy}
\end{equation}
The components are:
\begin{equation*}
\begin{aligned}
\partial_0\partial'_0y&=\partial_0\big(Ha'y-2aa'H^2\Delta x_0\big)=Ha'\big(Hay+2aa'H^2\Delta x_0\big)-2H^3a^2a'\Delta x_0+2aa'H^2\\
&=H^2aa'\big\{y-2aH\Delta x_0+2a'H\Delta x_0+2\big\},\\
\partial_0\partial'_jy&=H^2aa'\{-2aH\Delta x_j\},\qquad \partial_i\partial'_0y=H^2aa'\{2a'H\Delta x_i\},\qquad \partial_i\partial'_jy=H^2aa'\{-2\eta_{ij}\}.
\end{aligned}
\end{equation*}
The first line uses $\partial_0\Delta x_0=-1$. All components combine into a single expression:
\begin{resultbox}
\begin{equation}
\frac{\partial^2y}{\partial x^\rho\partial x'^\sigma}=H^2aa'\big\{y\,\delta^0_\rho\delta^0_\sigma-2a\,\delta^0_\rho H\Delta x_\sigma+2a'H\Delta x_\rho\delta^0_\sigma-2\eta_{\rho\sigma}\big\}.
\tag{I.71}\label{eq:ch05:I71}
\end{equation}
\end{resultbox}
\noindent This is Eq.~\eqI{71} of the paper. The lower-index convention $\Delta x_\sigma=\eta_{\sigma\nu}\Delta x^\nu$ is needed; the paper does not state it explicitly.

For the $\rho=\sigma=0$ component, use $y=aa'H^2\|\Delta\vec x\|^2-aa'H^2\Delta\eta^2$ and $-2H(a-a')\Delta x_0=2H^2aa'\Delta\eta^2$
(the latter from $H(a-a')=H^2aa'\Delta\eta$, Eq.~\eqref{eq:ch05:t12:aux} below). Then $y+2+2H^2aa'\Delta\eta^2=2-y+2aa'H^2\|\Delta\vec x\|^2$, and
\begin{equation}
\begin{aligned}
&\frac{\partial^2y}{\partial x^0\partial x'^0}=H^2aa'\big[2-y+2aa'H^2\|\Delta\vec x\|^2\big],\qquad
\frac{\partial^2y}{\partial x^0\partial x'^j}=H^2aa'\times-2aH\Delta x_j,\\
&\frac{\partial^2y}{\partial x^i\partial x'^0}=H^2aa'\times2a'H\Delta x_i,\qquad
\frac{\partial^2y}{\partial x^i\partial x'^j}=H^2aa'\times-2\eta_{ij}.
\end{aligned}
\tag{I.72--73}\label{eq:ch05:I72}
\end{equation}
These are \eqI{72}--\eqI{73}. They also agree with the table \eqref{eq:ch3-yder} of Task~8.
\eqref{eq:ch05:I72} is the regular part. The $\order{\delta}$ singular part of $\partial_0\partial'_0y$ was treated in Sections~\ref{sec:ch05:t11:69} and~\ref{sec:ch05:t11:delta}.

% -------------------------------------------------------------------
\subsection[Eq. (I.74)]{Eq.~(I.74) \tD}

For $F=I^2[f]$, Eqs.~\eqref{eq:ch05:I69} and \eqref{eq:ch05:t11:dy} give $\partial_iy\,\partial'_iy=-4(aa'H^2)^2(\Delta x_i)^2$ (no sum) and $\sum_i\partial_i\partial'_iy=-2(D-1)aa'H^2$. Hence
\begin{equation*}
\chVdd\,I^2[f]=f\sum_i\partial_iy\,\partial'_iy+I[f]\sum_i\partial_i\partial'_iy=-4(aa'H^2)^2\|\Delta\vec x\|^2f-2(D-1)aa'H^2\,I[f].
\end{equation*}
Divide by $-4aa'H^2$. Since $aa'$ is constant in space, it commutes with $\chVdd$:
\begin{resultbox}
\begin{equation}
aa'H^2\|\Delta\vec x\|^2f(y)=-\tfrac12(D-1)\,I[f](y)-\frac{\chVdd}{4aa'H^2}\,I^2[f](y).
\tag{I.74}\label{eq:ch05:I74}
\end{equation}
\end{resultbox}
\noindent This is \eqI{74}. The change $I^2[f]\to I^2[f]+c_1y$ produces $+\frac12(D-1)c_1$ on the right-hand side. This cancels against $I[f]\to I[f]+c_1$.

% -------------------------------------------------------------------
\subsection[Eqs. (I.75)--(I.78)]{Eqs.~(I.75)--(I.78) \tD}

All four are individual components of \eqref{eq:ch05:I70}, with $\partial\partial'y$ taken from \eqref{eq:ch05:I72}.

\paragraph{\eqI{75}, $\rho=\sigma=0$.}
Write $\partial_0\partial'_0y=H^2aa'(2-y)+2H^2aa'\cdot aa'H^2\|\Delta\vec x\|^2$.
The second piece multiplies $g\equiv I[f'A']$ and is converted by applying \eqref{eq:ch05:I74} to $g$:
\begin{equation*}
2H^2aa'\Big[-\tfrac12(D-1)I^2[f'A']-\frac{\chVdd}{4aa'H^2}I^3[f'A']\Big]=-(D-1)H^2aa'\,I^2[f'A']-\tfrac12\chVdd\,I^3[f'A'].
\end{equation*}
Therefore
\begin{resultbox}
\begin{equation}
f\,\partial_0\partial'_0A=\partial_0\partial'_0I^2[fA'']-\tfrac12\chVdd\,I^3[f'A']+H^2aa'\big\{(2-y)\,I[f'A']-(D-1)\,I^2[f'A']\big\}.
\tag{I.75}\label{eq:ch05:I75}
\end{equation}
\end{resultbox}

\paragraph{\eqI{76}, $\rho=0$, $\sigma=j$.}
Since $\partial'_jI^2[f'A']=I[f'A']\,\partial'_jy=-2aa'H^2\Delta x_j\,I[f'A']$, one has
\begin{equation*}
\partial_0\partial'_jy\;I[f'A']=-2H^3a^2a'\Delta x_j\,I[f'A']=Ha\,\partial'_jI^2[f'A'],
\end{equation*}
and therefore
\begin{equation}
f\,\partial_0\partial'_jA=\partial_0\partial'_jI^2[fA'']+Ha\,\partial'_jI^2[f'A'] .
\tag{I.76}\label{eq:ch05:I76}
\end{equation}

\paragraph{\eqI{77}, $\rho=i$, $\sigma=0$.}
In the same way, $\partial_iI^2[f'A']=2aa'H^2\Delta x_i\,I[f'A']$ gives $\partial_i\partial'_0y\;I[f'A']=Ha'\,\partial_iI^2[f'A']$:
\begin{equation}
f\,\partial_i\partial'_0A=\partial_i\partial'_0I^2[fA'']+Ha'\,\partial_iI^2[f'A'] .
\tag{I.77}\label{eq:ch05:I77}
\end{equation}

\paragraph{\eqI{78}, $\rho=i$, $\sigma=j$.}
\begin{equation}
f\,\partial_i\partial'_jA=\partial_i\partial'_jI^2[fA'']-2H^2aa'\eta_{ij}\,I[f'A'] .
\tag{I.78}\label{eq:ch05:I78}
\end{equation}

All four identities coincide with \eqI{75}--\eqI{78} of the paper. The footnote to \eqI{75} (naive derivative) is fully consistent with Section~\ref{sec:ch05:t11:delta}.

% -------------------------------------------------------------------
\subsection[Eqs. (I.79)--(I.80): normal part of the first term of (I.47)]{Eqs.~(I.79)--(I.80): the normal part of the first term of (I.47) \tD}

Take $f=A$ and the trace of \eqref{eq:ch05:I78}: $\sum_i\eta_{ii}=D-1$ and $f'A'=A'^2$. Thus
\begin{equation*}
A\,\chVdd A=\chVdd\,I^2[AA'']-2(D-1)H^2aa'\,I[A'^2].
\end{equation*}
The outer $\chVdd$ passes through $(aa')^{D-2}$, which is constant in space:
\begin{resultbox}
\begin{align}
\chVdd\big[(aa')^{D-2}A\,\chVdd A\big]
&=(aa')^{D-2}(\chVdd)^2\,I^2[AA'']-2(D-1)H^2(aa')^{D-1}\,\chVdd\,I[A'^2]
\tag{I.79}\label{eq:ch05:I79}\\
&=(aa')^{D-2}\nabla^4I^2[AA''](y)+2(D-1)H^2(aa')^{D-1}\nabla^2I[A'^2](y).
\tag{I.80}\label{eq:ch05:I80}
\end{align}
\end{resultbox}
\noindent The second line uses homogeneity: $\chVdd F(y)=-\nabla^2F(y)$, and $(\chVdd)^2=\nabla^4$.
\eqref{eq:ch05:I79} and \eqref{eq:ch05:I80} are \eqI{79} and \eqI{80} of the paper.
This is the square bracket without the factor $\kappa^2$.
In the re-derived \eqI{47} the coefficient of this term is $-\kappa^2$, while the paper prints $+\kappa^2$ (Task~7). This sign is carried over to Table~IV.

\subsubsection[The divergence claim after (I.80)]{The divergence claim after (I.80) \tC}
\label{sec:ch05:t11:power}

$y^s\sim\Delta x^{2s}$ is locally integrable in $D$ dimensions if and only if $s>-\frac D2$; at $D=4$ this means $s>-2$.
The powers in $A$ are $y^{1-D/2}$, $y^{2-D/2}$, a constant, and $y^n$. This gives:
\begin{center}
\small
\begin{tabular}{@{}lll@{}}
\toprule
Function & Powers & Not locally integrable at $D=4$\\
\midrule
$I^2[AA'']$ & $y^{2-D},\ y^{3-D},\ y^{4-D},\ y^{1-D/2},\ y^{2-D/2},\ y^{3-D/2},\ldots$ & only $y^{2-D}$ ($s=-2$, marginal)\\
$I[A'^2]$ & $y^{1-D},\ y^{2-D},\ y^{3-D},\ y^{1-D/2},\ y^{2-D/2},\ldots$ & $y^{1-D}$ ($s=-3$) and $y^{2-D}$ ($s=-2$)\\
\bottomrule
\end{tabular}
\end{center}
This confirms the paper's statement that only the first term (respectively the first two terms) can diverge.
\paragraph{A caveat.} In $I^2[AA'']$ the coefficients of the $y^{1-D/2}$ and $y^{3-D}$ terms each carry a pole ($c_f=k\pihat$ and $c_1$, respectively).
Combined, however, they are finite as $D\to4$, just like the $\ln y$ in \eqI{62}.
The statement ``one may set $D=4$ in the higher terms'' therefore means that these two terms are taken together, i.e.\ in the form of \eqI{62}.

% -------------------------------------------------------------------
\subsection{Comparison with the paper}
\label{sec:ch05:t11:comp}

\begin{center}
\small
\begin{tabular}{@{}l>{\raggedright\arraybackslash}p{0.42\textwidth}>{\raggedright\arraybackslash}p{0.17\textwidth}>{\raggedright\arraybackslash}p{0.16\textwidth}@{}}
\toprule
Equation & Derived here & Paper & Status\\
\midrule
\eqI{68} & definition \eqref{eq:ch05:I68}; constant ambiguity harmless & same & \chVyes\\
\eqI{69} & naive derivative: \eqref{eq:ch05:I69}; distributional derivative: extra $\gamma_F\delta^D$ term, \eqref{eq:ch05:t11:distr} & ``never a delta function'' & \chVyes{} (naive reading)\\
\eqI{70} & \eqref{eq:ch05:I70} & same & \chVyes\\
\eqI{71} & \eqref{eq:ch05:I71}, with lower-index $\Delta x_\sigma$ & same & \chVyes\\
\eqI{72}--\eqI{73} & \eqref{eq:ch05:I72} & same & \chVyes\\
\eqI{74} & \eqref{eq:ch05:I74} & same & \chVyes\\
\eqI{75}--\eqI{78} & \eqref{eq:ch05:I75}--\eqref{eq:ch05:I78} & same & \chVyes\\
\eqI{79}--\eqI{80} & \eqref{eq:ch05:I79}--\eqref{eq:ch05:I80} & same & \chVyes\\
divergence count & Section~\ref{sec:ch05:t11:power} & same & \chVyes\\
\bottomrule
\end{tabular}
\end{center}
No error in the paper was found in this part.
Since \eqI{68}--\eqI{80} do not depend on the sign convention, the overall sign question of \eqI{47} does not enter here.

\paragraph{Accumulated findings (Tasks 2--11).}
Identifiers in parentheses (S$n$, B$n$, C$n$) refer to the classified list of findings in Chapter~\ref{ch:summary}.
\begin{center}
\small
\begin{tabular}{@{}>{\raggedright\arraybackslash}p{0.30\textwidth}>{\raggedright\arraybackslash}p{0.33\textwidth}>{\raggedright\arraybackslash}p{0.30\textwidth}@{}}
\toprule
Location & Finding & Effect\\
\midrule
\eqI{7}, \eqI{14}, Table~II, \eqI{39}, the phrase ``4th term'' & typographical errors (C3--C7) & none\\
\eqI{40}, \eqI{59} & finite local errors (B1, B2) & absorbable into $\Delta\alpha_{2,3}$\\
\eqI{32}--\eqI{33}, \eqI{47}--\eqI{49} and beyond & 4-point and 3-point contributions reversed together (S1) & bare $M^2=-$(definition)\\
\eqI{60}--\eqI{67} & local cancellation against \eqI{44} verified & ---\\
\eqI{68}--\eqI{80} & agreement; naive-derivative qualification of \eqI{69} & \textbf{to be followed up:} is the $c_f\chVlead$ part differentiated naively in Tables~IV--VI?\\
\bottomrule
\end{tabular}
\end{center}

% -------------------------------------------------------------------
\subsection[Code checks]{Code checks \tC}
\label{sec:ch05:t11:code}

Script: \texttt{code/05\_normal\_terms\_tableIII/check\_68\_80.py} (SymPy; about 16--18\,s). It was actually run; its output is in \texttt{expected\_output/check\_68\_80.txt}.
It does not replace the analytic derivation above.
\paragraph{Method.} The functions are $f=y^\alpha$ and $A=y^\beta$ with generic exponents, so no integration constants arise. The script uses explicit conformal coordinates and the true $y(x;x')$.
The two sides are compared at random rational points ($\alpha=3/7$, $\beta=-5/3$) to 25 significant digits.
\begin{center}
\small
\begin{tabular}{@{}l>{\raggedright\arraybackslash}p{0.62\textwidth}l@{}}
\toprule
Check & What & Result\\
\midrule
(A) & \eqI{70}: all 16 components at $D=4$; also the consistency $I[fA'']=fA'-I[f'A']$ & \texttt{True}\\
(B) & \eqI{71} (all components, exact symbolic, lower-index $\Delta x_\sigma$) and the $(0,0)$ component of \eqI{72} & \texttt{True}\\
(C) & \eqI{74} at $D=4$ and $D=5$ & \texttt{True}\\
(D) & \eqI{75}--\eqI{78} (naive derivatives) at $D=4$ and $D=5$; \eqI{78} for both $i=j$ and $i\neq j$ & \texttt{True}\\
(E) & \eqI{79}--\eqI{80} at $D=4$ and $D=5$, with the left-hand side differentiated by brute force & \texttt{True}\\
(F) & \eqref{eq:ch05:t11:gamma}: with the model $f=c+d(4/y)^{D/2-1}+e(4/y)^{D/2-2}$ and $A=\mathcal N[c_0(4/y)^{D/2-1}+c_1(4/y)^{D/2-2}]$, the $y^{1-D/2}$ component of $I^2[fA'']$ is exactly $c\,\mathcal Nc_0\,4^{D/2-1}y^{1-D/2}=c_f\chVlead$ & \texttt{True}\\
(G) & power lists of Section~\ref{sec:ch05:t11:power}, generated with symbolic $D$, with the $D=4$ integrability criterion applied & as in the table\\
\bottomrule
\end{tabular}
\end{center}
Output of checks (F) and (G), lightly wrapped:
{\footnotesize
\begin{verbatim}
== (F) delta content: the y^{1-D/2} part of I^2[f A''] is c_f * A_lead
   (c_f = constant part of f) ==
  y^{1-D/2} part of I^2[fA''] : 2**(D - 2)*N*Y**(1 - D/2)*c*c0
     expected c*N*c0*4^(D/2-1)*Y^(1-D/2): True
  powers present in I^2[fA''] (exponent of Y): [1 - D/2, 2 - D, 2 - D/2, 3 - D, 4 - D]
== (G) power counting after (80), D=4: y^s is locally integrable in D dims iff s > -D/2 ==
  I^2[AA'']: [(2 - D, 'DIVERGENT'), (1 - D/2, 'finite'), (3 - D, 'finite'),
              (2 - D/2, 'finite'), (4 - D, 'finite'), (3 - D/2, 'finite')]
  I[A'^2]  : [(1 - D, 'DIVERGENT'), (2 - D, 'DIVERGENT'), (1 - D/2, 'finite'),
              (3 - D, 'finite'), (2 - D/2, 'finite'), (1, 'finite')]
\end{verbatim}
}
\noindent The rerun reproduces the original output in the Task~11 log. The only difference is the order of the two entries $y^{4-D}$ and $y^{2-D/2}$ in line (G).
Both have exponent $0$ at $D=4$, and their order was originally left to Python's set ordering; it is now fixed by a secondary sort key.

% -------------------------------------------------------------------
\subsection{Calculation record}
\begin{itemize}[leftmargin=14pt,itemsep=1pt]
\item \textbf{Derived:} \eqI{68} and its constant ambiguity; \eqI{69} in the naive form \eqref{eq:ch05:I69} and in the distributional form \eqref{eq:ch05:t11:distr};
\eqI{70} \eqref{eq:ch05:I70}, and the identity of its delta content with the local contribution \eqref{eq:ch05:t11:gamma};
\eqI{71}--\eqI{73} \eqref{eq:ch05:I71}--\eqref{eq:ch05:I72}; \eqI{74} \eqref{eq:ch05:I74}; \eqI{75}--\eqI{78} \eqref{eq:ch05:I75}--\eqref{eq:ch05:I78};
\eqI{79}--\eqI{80} \eqref{eq:ch05:I79}--\eqref{eq:ch05:I80}; the divergence count (Section~\ref{sec:ch05:t11:power}).
\item \textbf{Checks passed:} SymPy checks (A)--(G), actually run.
\item \textbf{Discrepancies with the paper:} none.
\item \textbf{To be followed up:} in the normal contributions $\partial\partial'I^2[fA'']$ must be taken naively. Taken as a distribution, it would count the local contribution of \eqI{54}--\eqI{58} twice (with a pole for $I=A$).
This is to be checked when Tables~IV--VI are derived.
\end{itemize}

% ===================================================================
\section[Task 12: Eqs. (I.81)--(I.93) and Table III]{Task 12: time-derivative identities, Table III and the covariant d'Alembertian, Eqs.~(I.81)--(I.93)}
\label{sec:ch05:t12}

\begin{discbox}[Summary of findings (Task 12)]
\begin{itemize}[leftmargin=12pt,itemsep=2pt]
\item \textbf{\eqI{81}--\eqI{88}, \eqI{90}, \eqI{91}--\eqI{92} and \eqI{93}} are identities that do not depend on the sign convention. All were derived and checked.
\item \textbf{\eqI{89}}: obtained by applying \eqI{88} to the nonlocal term of the re-derived \eqI{67}. The derived \eqI{89} is $-1$ times the printed one.
This difference is inherited from \eqI{47} (Tasks~7 and 10).
In addition, the second term of the printed \eqI{89} is \textbf{missing a factor $H^2$}: it should read $-4H^2(aa')^2(a+a')^2\nabla^2(yA')$ instead of $-4(aa')^2(a+a')^2\nabla^2(yA')$.
Dimensional analysis and a numerical check confirm this. Table~III gives this row correctly, with $H^2$; the error is only in the printing of \eqI{89}.
\item \textbf{Table~III}: the coefficients of all three rows were derived. The derived table is exactly $-1$ times the printed one (inherited sign). The row structure and the magnitudes agree exactly.
In the paper convention the derived table coincides with the printed one.
\item \textbf{Distributional remark.} In \eqI{91}--\eqI{92} the inner derivatives are the mixed $(0,j)$ components, which contain no delta function.
\eqI{92} therefore also holds in the distributional sense, and the $\square$ in it is a \emph{distributional} operator.
\eqI{93}, by contrast, is a pointwise identity: acting as a distribution on the $y^{1-D/2}$ component, $\square$ produces $\delta^D$, whereas \eqI{93} does not.
If $\square I^2[AA'']$ in \eqI{92} is evaluated pointwise with \eqI{93}, a legitimate local term $-ic_f\,a^{D-2}\nabla^2\delta^D$ (with a pole for $I=A$, $c_f=k\pihat$) is lost.
This is to be checked in Tables~IV--VI.
\end{itemize}
\end{discbox}

\paragraph{Notation specific to this section.}
$x\equiv y/4$ is used only in Table~III. $h\equiv H^2/16\pi^2$. $I[f]$ is the indefinite integral \eqref{eq:ch05:I68}.
\paragraph{Input principle.} Only results \emph{derived} in earlier tasks are used:
\eqI{67} from Task~10 (Eqs.~\eqref{eq:ch04:67D} and \eqref{eq:ch04:67}), \eqI{62} from Task~10 (Eq.~\eqref{eq:ch04:62}), and \eqI{74} from Task~11, Eq.~\eqref{eq:ch05:I74}.

% -------------------------------------------------------------------
\subsection{Goal, inputs and scope}

\paragraph{Targets \tP.}
The targets are \eqI{81}--\eqI{89} (p.~124005-8, left column), and \eqI{90}--\eqI{93} and Table~III (p.~124005-8, right column).
The printed equations were not used as input. They enter only in the comparison, Section~\ref{sec:ch05:t12:comp}.
Where the printed form differs from ours, it is retyped there.

\paragraph{Inputs.}
\begin{itemize}[leftmargin=14pt]
\item[(G1)] \tP\ \eqI{11}: $y=aa'H^2\{\|\Delta\vec x\|^2-\Delta\eta^2\}$ (regular part, $\delta\to0$).
\item[(G2)] \tD\ The re-derived \eqI{67} (Task~10, Eqs.~\eqref{eq:ch04:67D} and \eqref{eq:ch04:67}), in the \emph{definition sign}. Its nonlocal logarithm term in general $D$ is
\begin{equation}
\begin{aligned}
-iM^2_{\rm 3pt,log}&\supset\kappa^2k\,\frac{2}{D-3}\,H(aa')^{D-2}\big(a\partial_0+a'\partial'_0\big)\nabla^2A(y)\\
&\xrightarrow{D=4}\;\frac{\kappa^2H^2}{8\pi^2}\,2H(aa')^2\big(a\partial_0+a'\partial'_0\big)\nabla^2A(y).
\end{aligned}
\label{eq:ch05:t12:G2}
\end{equation}
In the paper's \eqI{67} this term appears as $-2H(\ldots)$ (paper convention).
\item[(G3)] \tD\ \eqI{62}: $A(y)=h\big\{4/y-2\ln(y/4)-1+\order{D-4}\big\}$.
\item[(G4)] \tD\ \eqI{74}: $aa'H^2\|\Delta\vec x\|^2f=-\tfrac12(D-1)I[f]-\dfrac{\chVdd}{4aa'H^2}I^2[f]$, Eq.~\eqref{eq:ch05:I74}.
\item[(G5)] \tP\ \eqI{42}: $a^D\square=\partial_\mu\big(a^{D-2}\eta^{\mu\nu}\partial_\nu\big)$; and \tD\ \eqI{76} (Task~11, Eq.~\eqref{eq:ch05:I76}).
\end{itemize}

\paragraph{Scope.}
\eqI{94}--\eqI{96} and Tables~IV--V are not treated in this task.
The identities are derived in general $D$. Table~III is given at $D=4$, as in the paper; the justification is in Task~10, Sec.~\ref{sec:ch04:Dto4}.
There, every term of the logarithm contribution is either $\delta^D$-local with a finite coefficient, or a smooth function times derivatives of a \emph{single} propagator.
Derivatives of a single propagator are well-defined distributions, analytic near $D=4$.

% -------------------------------------------------------------------
\subsection[Eq. (I.81)]{Eq.~(I.81) \tD}

$y$ depends on $\vec x$ and $\vec x'$ only through $\Delta\vec x$. Hence
\begin{equation}
\partial'_if(y)=f'(y)\,\partial'_iy=-f'(y)\,\partial_iy=-\partial_if(y).
\tag{I.81}\label{eq:ch05:I81}
\end{equation}

% -------------------------------------------------------------------
\subsection[Eqs. (I.82)--(I.83)]{Eqs.~(I.82)--(I.83) \tD}

\paragraph{An auxiliary identity.} Since $a=-1/(H\eta)$, we have $\frac1a-\frac1{a'}=-H\eta+H\eta'=-H\Delta\eta$. Hence
\begin{equation}
H\Delta\eta=\frac{a-a'}{aa'},\qquad a'H\Delta\eta=1-\frac{a'}a,\qquad aH\Delta\eta=\frac a{a'}-1 .
\label{eq:ch05:t12:aux}
\end{equation}
Using $\partial_0(aa'H^2)=Ha\cdot aa'H^2$ and $\partial_0\big(\|\Delta\vec x\|^2-\Delta\eta^2\big)=-2\Delta\eta$:
\begin{equation}
\frac{\partial y}{\partial x^0}=Ha\,y-2aa'H^2\Delta\eta=Ha\big(y-2a'H\Delta\eta\big)\overset{\eqref{eq:ch05:t12:aux}}{=}Ha\Big(y-2+2\frac{a'}a\Big).
\tag{I.82}\label{eq:ch05:I82}
\end{equation}
Similarly, with $\partial'_0(\ldots)=+2\Delta\eta$:
\begin{equation}
\frac{\partial y}{\partial x'^0}=Ha'\,y+2aa'H^2\Delta\eta=Ha'\big(y+2aH\Delta\eta\big)\overset{\eqref{eq:ch05:t12:aux}}{=}Ha'\Big(y-2+2\frac a{a'}\Big).
\tag{I.83}\label{eq:ch05:I83}
\end{equation}
\eqref{eq:ch05:I82} and \eqref{eq:ch05:I83} are \eqI{82}--\eqI{83} of the paper.

% -------------------------------------------------------------------
\subsection[Eqs. (I.84)--(I.88)]{Eqs.~(I.84)--(I.88) \tD}

\paragraph{\eqI{84}.}
\begin{equation}
\begin{aligned}
(\partial_0+\partial'_0)f&=f'\big[\partial_0y+\partial'_0y\big]=f'\big[Hay-2aa'H^2\Delta\eta+Ha'y+2aa'H^2\Delta\eta\big]\\
&=H(a+a')\,y\,f'(y).
\end{aligned}
\tag{I.84}\label{eq:ch05:I84}
\end{equation}

\paragraph{\eqI{85}.}
\begin{equation*}
(a'\partial_0+a\partial'_0)f=f'\big[a'(Hay-2aa'H^2\Delta\eta)+a(Ha'y+2aa'H^2\Delta\eta)\big]=f'\,Haa'\big[2y+2H\Delta\eta\,(a-a')\big].
\end{equation*}
By \eqref{eq:ch05:t12:aux}, $H(a-a')=H^2aa'\Delta\eta$, so $2H\Delta\eta(a-a')=2H^2aa'\Delta\eta^2$.
The square bracket becomes $2aa'H^2\big(\|\Delta\vec x\|^2-\Delta\eta^2\big)+2aa'H^2\Delta\eta^2=2aa'H^2\|\Delta\vec x\|^2$:
\begin{equation}
(a'\partial_0+a\partial'_0)f(y)=2Haa'\times aa'H^2\|\Delta\vec x\|^2f'(y).
\tag{I.85}\label{eq:ch05:I85}
\end{equation}

\paragraph{\eqI{86}.}
Apply (G4) with $f\to f'$, writing $I[f']=f+c$ and $I^2[f']=I[f]+cy+c'$:
\begin{equation*}
aa'H^2\|\Delta\vec x\|^2f'=-\tfrac12(D-1)(f+c)-\frac{\chVdd}{4aa'H^2}\big(I[f]+cy\big)=-\tfrac12(D-1)f-\frac{\chVdd}{4aa'H^2}I[f].
\end{equation*}
The $c$ terms cancel because $\chVdd y=\sum_i\partial_i\partial'_iy=-2(D-1)aa'H^2$.
Insert this into \eqref{eq:ch05:I85} and use $\chVdd F(y)=-\nabla^2F(y)$:
\begin{resultbox}
\begin{equation}
(a'\partial_0+a\partial'_0)f(y)=-(D-1)Haa'f(y)+\frac{\nabla^2}{2H}\,I[f](y).
\tag{I.86}\label{eq:ch05:I86}
\end{equation}
\end{resultbox}

\paragraph{\eqI{87}--\eqI{88}.}
$(a+a')(\partial_0+\partial'_0)-(a'\partial_0+a\partial'_0)=a\partial_0+a'\partial'_0$ is an operator identity; this is \eqI{87}. With \eqref{eq:ch05:I84} and \eqref{eq:ch05:I86}:
\begin{resultbox}
\begin{equation}
(a\partial_0+a'\partial'_0)f(y)=H(a+a')^2y\,f'(y)+(D-1)Haa'f(y)-\frac{\nabla^2}{2H}\,I[f](y).
\tag{I.88}\label{eq:ch05:I88}
\end{equation}
\end{resultbox}
\noindent \eqref{eq:ch05:I84}--\eqref{eq:ch05:I88} are \eqI{84}--\eqI{88} of the paper. All hold in general $D$ and do not depend on the integration constant.

% -------------------------------------------------------------------
\subsection[Eq. (I.89)]{Eq.~(I.89) \tD}
\label{sec:ch05:t12:89}

Since $a$ and $a'$ depend only on time, $(a\partial_0+a'\partial'_0)$ commutes with $\nabla^2$.
Apply \eqref{eq:ch05:I88} with $f=A$ to (G2). In general $D$ (definition sign):
\begin{equation}
\kappa^2k\,\frac{2}{D-3}\,H(aa')^{D-2}\,\nabla^2\Big[H(a+a')^2yA'+(D-1)Haa'A-\frac{\nabla^2}{2H}I[A]\Big].
\label{eq:ch05:t12:89D}
\end{equation}
At $D=4$ ($k\to2h=H^2/8\pi^2$, $\frac2{D-3}\to2$, $D-1\to3$):
\begin{equation*}
\frac{\kappa^2H^2}{8\pi^2}\Big\{2H^2(aa')^2(a+a')^2\nabla^2(yA')+6H^2(aa')^3\nabla^2A-(aa')^2\nabla^4I[A]\Big\},
\end{equation*}
that is,
\begin{resultbox}[Derived (I.89), definition sign]
\begin{equation}
\begin{aligned}
&\frac{\kappa^2H^2}{8\pi^2}\,2H(aa')^2\big(a\partial_0+a'\partial'_0\big)\nabla^2A\\
&\qquad=\frac{\kappa^2H^2}{16\pi^2}\Big\{+12(aa')^3H^2\nabla^2A+4H^2(aa')^2(a+a')^2\nabla^2(yA')-2(aa')^2\nabla^4I[A]\Big\}.
\end{aligned}
\tag{I.89$'$}\label{eq:ch05:I89}
\end{equation}
\end{resultbox}

\begin{discbox}[Comparison with the printed (I.89)]
\PI{} prints (paper convention):
\begin{equation*}
\begin{aligned}
&\frac{\kappa^2H^2}{8\pi^2}\,\big(-2H\big)(aa')^2\big(a\partial_0+a'\partial'_0\big)\nabla^2A(y)\\
&\qquad=\frac{\kappa^2H^2}{16\pi^2}\Big\{-12(aa')^3H^2\nabla^2A-4(aa')^2(a+a')^2\nabla^2(yA')+2(aa')^2\nabla^4I[A]\Big\}.
\end{aligned}
\end{equation*}
\begin{itemize}[leftmargin=12pt]
\item \textbf{Overall sign.} The overall sign is reversed. The left-hand side of the paper is $-2H(\ldots)$; this is the sign of \eqI{67} inherited from \eqI{47}.
In the paper convention our result coincides with the printed coefficients $-12$, $-4H^2$, $+2$.
\item \textbf{Missing $H^2$ in the second term.} In conformal coordinates $x^\mu$ has dimensions of length and $H$ of inverse length; $a$ and $y$ are dimensionless.
The terms $H^2\nabla^2A$ and $\nabla^4I[A]$ have dimension $L^{-4}A$, whereas the printed $(a+a')^2\nabla^2(yA')$ has dimension $L^{-2}A$.
The correct term is $-4H^2(aa')^2(a+a')^2\nabla^2(yA')$, and the second row of Table~III does contain $H^2$.
Numerical check (Section~\ref{sec:ch05:t12:code}, item (f)): with $H^2$ the identity \eqI{89} holds at every point tested; without $H^2$ it fails for $H\neq1$.
\end{itemize}
\end{discbox}

% -------------------------------------------------------------------
\subsection[Table III]{Table III \tD}
\label{sec:ch05:t12:tab3}

\paragraph{Normalization.}
The prefactor of \eqref{eq:ch05:I89} is $\frac{\kappa^2H^2}{16\pi^2}$. Write $A=h\cdot\chVAh$ with $\chVAh\equiv4/y-2\ln(y/4)-1$. Then
\begin{equation*}
\frac{\kappa^2H^2}{16\pi^2}\cdot h=\frac{\kappa^2H^4}{256\pi^4}=\frac{\kappa^2H^4}{(4\pi)^4}.
\end{equation*}
So the coefficients of Table~III are read off from the brace of \eqref{eq:ch05:I89} with $A\to\chVAh$.

\paragraph{The three rows.}
Terms are dropped as follows. Constants vanish under $\nabla^2$ or $\nabla^4$.
Moreover, $\nabla^2y=2(D-1)aa'H^2$ is constant in space, so $\nabla^4y=0$. This is what the paper means by ``terms which eventually vanish such as $\nabla^4y$''.
\begin{itemize}[leftmargin=14pt]
\item \textbf{Row 1}, $(aa')^3H^2\nabla^2$: $12\chVAh=12\big(\tfrac1x-2\ln x-1\big)\to\dfrac{12}{x}-24\ln x$. The constant $-12$ is dropped; $4/y=1/x$ and $\ln(y/4)=\ln x$.
\item \textbf{Row 2}, $(aa')^2(a+a')^2H^2\nabla^2$: $\chVAh'=-\tfrac4{y^2}-\tfrac2y$, so $y\chVAh'=-\tfrac4y-2$.
Hence $4y\chVAh'=-\tfrac{16}{y}-8\to-\dfrac4x$.
\item \textbf{Row 3}, $(aa')^2\nabla^4$: $I[\chVAh]=4\ln y-2y\ln\tfrac y4+y$, using $\int\ln\tfrac y4\,dy=y\ln\tfrac y4-y$.
Then $4\ln y=4\ln x+4\ln4$, and with $y=4x$ one finds $-2I[\chVAh]=-8\ln x-8\ln4+16x\ln x-8x\to-8\ln x+16x\ln x$. The constant and the term linear in $y$ are dropped.
\end{itemize}

\begin{center}
\small
\begin{tabular}{@{}llll@{}}
\toprule
External operator & Derived, definition sign & Derived, paper conv. & \PI{} prints (Table~III)\\
& $\big(\times\frac{\kappa^2H^4}{(4\pi)^4}\big)$ & $\big(\times\frac{\kappa^2H^4}{(4\pi)^4}\big)$ & $\big(\times\frac{\kappa^2H^4}{(4\pi)^4}\big)$\\
\midrule
$(aa')^3H^2\nabla^2$ & $+\frac{12}{x}-24\ln x$ & $-\frac{12}{x}+24\ln x$ & $-\frac{12}{x}+24\ln x$\\[2pt]
$(aa')^2(a+a')^2H^2\nabla^2$ & $-\frac4x$ & $\frac4x$ & $\frac4x$\\[2pt]
$(aa')^2\nabla^4$ & $-8\ln x+16x\ln x$ & $8\ln x-16x\ln x$ & $8\ln x-16x\ln x$\\
\bottomrule
\end{tabular}
\end{center}
Relative to the printed table, every derived entry has the opposite sign (definition sign versus paper convention). This is the inherited overall sign, not a new error.
The same computation starting from the paper's own \eqI{89}, with $H^2$ restored, reproduces the paper's Table~III exactly.
The table is thus the correct processing of the paper's own \eqI{89}; the only difference is the overall sign inherited from \eqI{47}.

\paragraph{Meaning of the external operators.}
$1/x$, $\ln x$ and $x\ln x$ are locally integrable at $D=4$. The external operators act as distributions, i.e.\ by integration by parts onto the test function.
At fixed $\delta>0$ these functions are smooth in space. So $\nabla^2$ and $\nabla^4$, which contain only spatial derivatives, produce no $\delta^D$ term.
Table~III has no local part, and the word ``nonlocal'' in its caption is correct in this sense.

% -------------------------------------------------------------------
\subsection[Eq. (I.90)]{Eq.~(I.90) \tD}

From \eqI{42}, $a^D\square=-\partial_0\big(a^{D-2}\partial_0\big)+a^{D-2}\nabla^2$. With $\partial_0a^{D-2}=(D-2)Ha^{D-1}$:
\begin{equation}
\square=-\frac1{a^2}\partial_0^2-\frac{(D-2)H}{a}\partial_0+\frac1{a^2}\nabla^2 .
\tag{I.90}\label{eq:ch05:I90}
\end{equation}

% -------------------------------------------------------------------
\subsection[Eqs. (I.91)--(I.92)]{Eqs.~(I.91)--(I.92) \tD}
\label{sec:ch05:t12:91}

\paragraph{\eqI{91}.} Use \eqref{eq:ch05:I76} with $f=A$, summed over $j$ ($f'A'=A'^2$):
\begin{equation*}
A\,\partial_0\partial'_jA=\partial_0\partial'_jI^2[AA'']+Ha\,\partial'_jI^2[A'^2].
\end{equation*}
$\partial'_j$ commutes with the time-dependent factor $(aa')^{D-2}$ and with $\partial_0$. It can therefore be moved outside, where it combines with the inner $\partial'_j$ into $\nabla'^2$:
\begin{equation}
\begin{aligned}
\partial_0\partial'_i\big[(aa')^{D-2}A\,\partial_0\partial'_iA\big]={}&\nabla'^2\partial_0\big[(aa')^{D-2}\partial_0I^2[AA'']\big]\\
&+H\nabla'^2\partial_0\big[a^{D-1}a'^{D-2}I^2[A'^2]\big].
\end{aligned}
\tag{I.91}\label{eq:ch05:I91}
\end{equation}

\paragraph{\eqI{92}.} By homogeneity, $\nabla'^2=\nabla^2$. Define $G\equiv I^2[AA'']$ and $K\equiv I^2[A'^2]$.
\begin{itemize}[leftmargin=14pt]
\item In the first term, \eqI{42} gives $\partial_0\big(a^{D-2}\partial_0G\big)=-a^D\square G+a^{D-2}\nabla^2G$.
\item In the second term, $\partial_0\big(a^{D-1}K\big)=(D-1)Ha^DK+a^{D-1}\partial_0K$.
\end{itemize}
\begin{resultbox}
\begin{equation}
\begin{aligned}
={}&-a^Da'^{D-2}\nabla^2\square I^2[AA'']+(aa')^{D-2}\nabla^4I^2[AA'']\\
&+Ha^{D-1}a'^{D-2}\nabla^2\partial_0I^2[A'^2]+(D-1)H^2a^Da'^{D-2}\nabla^2I^2[A'^2].
\end{aligned}
\tag{I.92}\label{eq:ch05:I92}
\end{equation}
\end{resultbox}
\noindent \eqref{eq:ch05:I91} and \eqref{eq:ch05:I92} are \eqI{91}--\eqI{92} of the paper.

\subsubsection[Distributional meaning]{Distributional meaning \tA}
\begin{itemize}[leftmargin=14pt]
\item The left-hand side of \eqref{eq:ch05:I91} is the normal part of the fourth term of \eqI{47}.
Its inner pair of derivatives, $\partial_0\partial'_i$, is mixed and contains no delta function even for the full $i\Delta_A$ (Task~8).
So the naive and distributional readings coincide, and \eqref{eq:ch05:I76} holds here in both senses.
\item The outer derivatives are always distributional. Distributional derivatives commute, so the step \eqref{eq:ch05:I91}$\to$\eqref{eq:ch05:I92} is exact in the distributional sense.
The $\square$ in \eqref{eq:ch05:I92} is a \textbf{distributional} operator.
\item $G=I^2[AA'']$ contains the component $c_f\chVlead$ (Task~11, Eq.~\eqref{eq:ch05:t11:gamma}). Here $c_f=k\pihat$, which has a pole.
From the propagator equation, $a^D\square\chVlead\supset i\delta^D$ (Task~8, route B, Eq.~\eqref{eq:ch3-routeB2}). Hence
\begin{equation}
-a^Da'^{D-2}\nabla^2\square I^2[AA'']\supset-i\,c_f\,a^{D-2}\nabla^2\delta^D(x-x')=-ik\pihat\,a^{D-2}\nabla^2\delta^D .
\label{eq:ch05:t12:loc}
\end{equation}
This term is legitimate. It is the exact counterpart of the local term $-ia^{D-2}kL\nabla^2\delta^D$ of $\mathcal V_4$ in Task~10, where the outer function is $kL$ instead of $c_f$.
\eqI{54}--\eqI{59} contain only the inner $(0,0)$ terms, so \eqref{eq:ch05:t12:loc} was not counted there.
\end{itemize}

% -------------------------------------------------------------------
\subsection[Eq. (I.93)]{Eq.~(I.93) \tD}
\label{sec:ch05:t12:93}

By \eqI{42}, $\square f(y)=f''(y)\,a^{-2}\eta^{\mu\nu}\partial_\mu y\,\partial_\nu y+f'(y)\,\square y$. Two quantities are needed.
\begin{itemize}[leftmargin=14pt]
\item With \eqref{eq:ch05:I82} and $\partial_iy=2aa'H^2\Delta x^i$:
\begin{equation*}
\begin{aligned}
\eta^{\mu\nu}\partial_\mu y\,\partial_\nu y&=-\big(Hay-2aa'H^2\Delta\eta\big)^2+4a^2a'^2H^4\|\Delta\vec x\|^2\\
&=-H^2a^2y^2+4H^3a^2a'y\Delta\eta+4aa'H^2y .
\end{aligned}
\end{equation*}
(The last step uses $4a^2a'^2H^4(\|\Delta\vec x\|^2-\Delta\eta^2)=4aa'H^2y$.)
Divide by $a^2$ and use $a'H\Delta\eta=1-a'/a$ from \eqref{eq:ch05:t12:aux}:
$-H^2y^2+4H^2y\,(1-a'/a)+4H^2y\,a'/a=H^2(4y-y^2)$.
\item $\nabla^2y=2(D-1)aa'H^2$ and $\partial_0^2y=2H^2a^2y-4H^3a^2a'\Delta\eta-2aa'H^2$. With \eqref{eq:ch05:I90}:
\begin{equation*}
\square y=-DH^2y+2DH^2\big(a'H\Delta\eta+a'/a\big)=DH^2(2-y).
\end{equation*}
(The $y$ terms are $-2H^2y-(D-2)H^2y$; the $\Delta\eta$ terms are $[4+2(D-2)]H^3a'\Delta\eta$; the $a'/a$ terms are $[2+2(D-1)]H^2a'/a$.)
\end{itemize}
\begin{resultbox}
\begin{equation}
\square f(y)=H^2\big[(4y-y^2)f''(y)+D(2-y)f'(y)\big]=\square'f(y).
\tag{I.93}\label{eq:ch05:I93}
\end{equation}
\end{resultbox}
\noindent The equality with $\square'$ follows from the symmetry $y(x;x')=y(x';x)$. This is \eqI{93} of the paper.

\paragraph{Remark.} \eqref{eq:ch05:I93} is pointwise.
The Task~12 report states that for $f=y^{1-D/2}$ the right-hand side vanishes identically (indicial equation), whereas as a distribution $\square y^{1-D/2}$ contains $\delta^D$.
\begin{notebox}[Editorial note added while preparing these notes]
Strictly, only the leading (indicial) power cancels for $f=y^{1-D/2}$, namely the coefficient of $y^{-D/2}$, which is $4s(s-1)+2Ds=0$ for $s=1-\frac D2$.
The next power survives. Pointwise,
$\square\,y^{1-D/2}=\tfrac{D(D-2)}{4}\,H^2\,y^{1-D/2}$, a locally integrable function with no delta function (we confirmed this with a one-line SymPy check).
The conclusion of the remark is unaffected: the pointwise formula \eqref{eq:ch05:I93} never produces $\delta^D$, while the distributional $\square$ acting on $y^{1-D/2}$ does.
\end{notebox}
\noindent Therefore, if $\square I^2[AA'']$ in \eqref{eq:ch05:I92} is evaluated pointwise with \eqref{eq:ch05:I93}, the legitimate local term \eqref{eq:ch05:t12:loc} is lost.
Conversely, the inner $(0,0)$ term in \eqref{eq:ch05:I75} was defined naively, and its local part was counted separately in \eqI{54}.
If it combines there with outer derivatives into a distributional $\square$, an extra delta function appears.
The consistent rule is: the outer operators are distributional, and $c_f\,\frac{i}{a^{D-2}}\delta^D$ is subtracted from the inner $(0,0)$ derivative, because it was already counted in \eqI{54}--\eqI{58}.
Whether this rule is respected in Tables~IV--VI is to be checked there.

% -------------------------------------------------------------------
\subsection{Comparison with the paper}
\label{sec:ch05:t12:comp}

\begin{center}
\small
\begin{tabular}{@{}>{\raggedright\arraybackslash}p{0.20\textwidth}>{\raggedright\arraybackslash}p{0.30\textwidth}>{\raggedright\arraybackslash}p{0.22\textwidth}>{\raggedright\arraybackslash}p{0.20\textwidth}@{}}
\toprule
Item & Derived here & Paper & Status\\
\midrule
\eqI{81} & \eqref{eq:ch05:I81} & same & \chVyes\\
\eqI{82}--\eqI{83} & \eqref{eq:ch05:I82}--\eqref{eq:ch05:I83} & same & \chVyes\\
\eqI{84}--\eqI{88} & \eqref{eq:ch05:I84}--\eqref{eq:ch05:I88}, general $D$ & same & \chVyes\\
\eqI{89}: overall sign & $+12,\ +4H^2,\ -2$ (definition sign) & $-12,\ -4,\ +2$ & \chVno{} (sign, inherited)\\
\eqI{89}: $H^2$ in 2nd term & present & absent & \chVno{} (typo)\\
Table~III & Section~\ref{sec:ch05:t12:tab3} & all three rows with opposite sign & \chVno{} (sign, inherited)\\
\eqI{90} & \eqref{eq:ch05:I90} & same & \chVyes\\
\eqI{91}--\eqI{92} & \eqref{eq:ch05:I91}--\eqref{eq:ch05:I92} & same & \chVyes\\
\eqI{93} & \eqref{eq:ch05:I93}, pointwise & same & \chVyes\\
\bottomrule
\end{tabular}
\end{center}

\paragraph{Accumulated findings (Tasks 2--12).}
\begin{center}
\small
\begin{tabular}{@{}>{\raggedright\arraybackslash}p{0.30\textwidth}>{\raggedright\arraybackslash}p{0.36\textwidth}>{\raggedright\arraybackslash}p{0.27\textwidth}@{}}
\toprule
Location & Finding & Effect\\
\midrule
\eqI{7}, \eqI{14}, Table~II, \eqI{39}, the phrase ``4th term'' & typographical errors (C3--C7) & none\\
\textbf{\eqI{89}} & $H^2$ missing in the second term (C8) & none (Table~III is correct)\\
\eqI{40}, \eqI{59} & finite local errors (B1, B2) & absorbable into $\Delta\alpha_{2,3}$\\
\eqI{32}--\eqI{33}, \eqI{47}--\eqI{49} and beyond & 4-point and 3-point contributions reversed together, Table~III included (S1) & bare $M^2=-$(definition)\\
\eqI{60}--\eqI{67} & local cancellation against \eqI{44} verified & ---\\
\eqI{68}--\eqI{80}, \eqI{81}--\eqI{93} & agreement & ---\\
\eqI{92}--\eqI{93} & is $\square$ distributional or pointwise? & \textbf{to be followed up:} is the local term \eqref{eq:ch05:t12:loc} present in Tables~IV--VI?\\
\bottomrule
\end{tabular}
\end{center}

% -------------------------------------------------------------------
\subsection[Code checks]{Code checks \tC}
\label{sec:ch05:t12:code}

Script: \texttt{code/05\_normal\_terms\_tableIII/check\_81\_93.py} (SymPy; about 20\,s). It was actually run; its output is in \texttt{expected\_output/check\_81\_93.txt}.
It does not replace the analytic derivation. The script uses explicit coordinates and the true $y(x;x')$.
\begin{center}
\small
\begin{tabular}{@{}l>{\raggedright\arraybackslash}p{0.66\textwidth}l@{}}
\toprule
Check & What & Result\\
\midrule
(a) & \eqI{81}--\eqI{83}, exact symbolic & \texttt{True}\\
(b) & \eqI{84}, \eqI{85}, \eqI{87} symbolic with a generic function $f(y)$; \eqI{86}, \eqI{88} with $f=y^\alpha$ at a random point, 25 digits & \texttt{True}\\
(c) & \eqI{90} in general $D$, from the definition \eqI{42}, on a test function $g(\eta,\vec x)$ & \texttt{True}\\
(d) & \eqI{91}--\eqI{92} with $A=y^\beta$ at $D=4,5$; left side by brute-force differentiation, $\square$ applied in coordinates via \eqI{42} & \texttt{True}\\
(e) & \eqI{93}: $\square$ and $\square'$ at $D=4,5$ & \texttt{True}\\
(f) & \eqI{89} and Table~III at $D=4$ with the true $A(y)$ of \eqI{62}, at three test points with $H\neq1$ & see below\\
\bottomrule
\end{tabular}
\end{center}
Results of check (f):
\begin{itemize}[leftmargin=14pt]
\item The derived nonlocal term of \eqI{67} plus the printed Table~III gives exactly $0$ at all three points. So derived $=-$paper.
\item \eqI{89} with $H^2$ restored equals $-1$ times the derived term at every point. The printed \eqI{89} (without $H^2$) does not at any point.
\item The sources of the table coefficients ($-12\chVAh$, $-4y\chVAh'$, $2I[\chVAh]$, paper convention) were expanded separately.
\end{itemize}
Output of block (f):
{\footnotesize
\begin{verbatim}
== (89) and Table III: pointwise, D=4, actual A(y) of (62) ==
  point 0: derived(67)_nonlocal + Table III(paper) = 0.0
           | derived - Table III(paper) = 0.0000509  (scale 0.0000254)
           (89) with H^2 restored == -derived : True
           (89) as printed (no H^2) == -derived : False
  point 1: derived(67)_nonlocal + Table III(paper) = 0.0
           | derived - Table III(paper) = 0.000202  (scale 0.000101)
           (89) with H^2 restored == -derived : True
           (89) as printed (no H^2) == -derived : False
  point 2: derived(67)_nonlocal + Table III(paper) = 0.0
           | derived - Table III(paper) = 0.000303  (scale 0.000152)
           (89) with H^2 restored == -derived : True
           (89) as printed (no H^2) == -derived : False
  Table III coefficients from A(y) (dropping terms killed by Lap / Lap^2):
    -12*A/h        = 24*log(x) + 12 - 12/x
    -4*y A'/h      = 8 + 4/x
    2*I[A]/h       = -16*x*log(x) + 8*x + 8*log(x) + 16*log(2)
\end{verbatim}
}
\paragraph{A note on the test points.}
In the first run recorded in the Task~12 log, $H$ was drawn as a random integer in $[5,20]$ divided by 10. It came out as $H=1$ at points 0 and 2.
There the printed \eqI{89} (without $H^2$) passed by accident (\texttt{True}).
The script was then changed to draw $H$ from $\{7/10,13/10,3/2,17/10\}$. The output above is from this final version, and our rerun reproduces it line by line.
At fixed $(\eta,\eta',\vec x,\vec x')$ the derived term and Table~III do not depend on $H$, which is why the printed scales are the same in both runs.

% -------------------------------------------------------------------
\subsection{Calculation record}
\begin{itemize}[leftmargin=14pt,itemsep=1pt]
\item \textbf{Derived:} \eqI{81} \eqref{eq:ch05:I81}; \eqI{82}--\eqI{83} \eqref{eq:ch05:I82}--\eqref{eq:ch05:I83};
\eqI{84}--\eqI{88} \eqref{eq:ch05:I84}--\eqref{eq:ch05:I88}, in general $D$ and independent of the integration constant;
\eqI{89} in general $D$ \eqref{eq:ch05:t12:89D} and at $D=4$ \eqref{eq:ch05:I89}; Table~III (Section~\ref{sec:ch05:t12:tab3});
\eqI{90} \eqref{eq:ch05:I90}; \eqI{91}--\eqI{92} \eqref{eq:ch05:I91}--\eqref{eq:ch05:I92} and their distributional meaning \eqref{eq:ch05:t12:loc}; \eqI{93} \eqref{eq:ch05:I93}.
\item \textbf{Checks passed:} SymPy checks (a)--(f), actually run.
\item \textbf{Discrepancies with the paper:}
the overall sign of \eqI{89} and Table~III is reversed. This is inherited from \eqI{47} and is not a new error.
The second term of \eqI{89} is missing $H^2$. This is a typographical error; Table~III is correct.
\item \textbf{To be followed up:} the $\square$ in \eqI{92} is distributional, while \eqI{93} is pointwise.
It must be checked whether the local term \eqref{eq:ch05:t12:loc}, which has a pole, is present in Tables~IV--VI.
\end{itemize}

% -------------------------------------------------------------------
\begin{notebox}[Forward note (results of later tasks, not part of Tasks 11--12)]
\begin{itemize}[leftmargin=12pt,itemsep=1pt]
\item \textbf{Naive inner derivatives.} Task~13 adopts naive inner derivatives in the normal contributions, following the paper's appendix text and the footnote to \eqI{75}.
Task~16 (Sec.~\ref{sec:ch08-distr}) finds that the possible double counting identified here and in Task~13 has exactly the form $c\,(a^D\square-a^{D-2}\nabla^2)\delta^D$, with $c=c(D)$ a constant and no $\ln a$.
This is the form of the counterterms, so it is absorbed completely into $\alpha_2$ and $\alpha_3$ and does not change the renormalized result.
The general question remains open beyond pointwise checks: does the external-operator reduction of the appendix produce no \emph{other} local differences? Task~16 records this point as only partially closed.
\item \textbf{Table~III in Table~X.} Task~17 rebuilds Table~X from Tables~III, V and VII and Eq.~\eqI{159}.
It finds that the paper added the first two rows of Table~III to the corresponding rows of Table~X, but omitted the third row, $8\ln x-16x\ln x$ (paper convention), from the $(aa')^2\nabla^4$ row.
The corrected row is $\frac{34}{3}\ln x-40x\ln x+24x^2\ln x-36x^2$, where the paper prints $\frac{10}{3}\ln x-24x\ln x+24x^2\ln x-36x^2$.
This is discussed in the chapter on renormalization and Table~X.
\end{itemize}
\end{notebox}
% (hypertexnames is left false: notes.tex sets it globally, and switching it back on here
%  would create duplicate equation anchors in the following chapters.)
